% latexmk -pdf camera_pose_covariance.tex && latexmk -c camera_pose_covariance.tex

\documentclass{article}
\usepackage{amsmath}
\usepackage{amssymb}
\usepackage{mathtools}
\usepackage{bm}
\usepackage{hyperref}

\newcommand{\tran}{\top}

\newcommand{\B}{\mathrm{B}}
\newcommand{\C}{\mathrm{C}}
\newcommand{\W}{\mathrm{W}}

\newcommand{\inW}{{}^{\W}\!}
\newcommand{\inB}{{}^{\B}\!}
\newcommand{\inC}{{}^{\C}\!}
\newcommand{\inCs}{{}^{\C}}

\newcommand{\Ab}{\bm{A}}
\newcommand{\Hb}{\bm{H}}
\newcommand{\Ib}{\bm{I}}
\newcommand{\Rb}{\bm{R}}
\newcommand{\ob}{\bm{o}}
\newcommand{\Sb}{\bm{\Sigma}}
\newcommand{\J}{\bm{J}}
\newcommand{\m}{\bm{m}}
\newcommand{\xib}{\bm{\xi}}
\newcommand{\pb}{\bm{p}}
\newcommand{\nub}{\bm{\nu}}
\newcommand{\omegab}{\bm{\omega}}
\newcommand{\db}{\bm{d}}
\newcommand{\vb}{\bm{v}}
\newcommand{\wb}{\bm{w}}
\newcommand{\Zb}{\bm{0}}

\newcommand{\Rbt}{\bm{\tilde{R}}}
\newcommand{\pbt}{\bm{\tilde{p}}}
\newcommand{\fb}{\bm{f}}
\newcommand{\rb}{\bm{r}}
\newcommand{\Pb}{\bm{P}}
\newcommand{\epsb}{\bm{\varepsilon}}


\begin{document}
This note derives the camera pose covariance $\inCs \Sb_\C$ of \href{pose_estimation.pdf}{\texttt{pose\_estimation.pdf}} and the estimate of the pixel noise it is scaled with.
For $N$ detected marker corners, the image Jacobian $\J$ stacks the $2 \times 6$ Jacobians of all corners into a $2N \times 6$ matrix, so each row belongs to one pixel coordinate, $u$ or $v$, of one corner.

\paragraph{Covariance of the fitted pose.}
Let $\xib$ be the true camera pose and $\fb(\xib)$ the projection of the marker points into the image at that pose, which maps a camera pose to the stacked $2N$ pixel coordinates of the marker corners. We model the detected corner coordinates $\m \in \mathbb{R}^{2N}$ as this projection plus independent pixel noise of equal variance $\sigma^2$ on every coordinate,
\[
  \m = \fb(\xib) + \epsb, \qquad \epsb \sim \mathcal{N}(\Zb, \sigma^2 \Ib_{2N}).
\]
SQPnP returns the pose $\hat\xib$ that minimizes the reprojection cost $E = \frac{1}{2} \lVert \rb \rVert^2$ of the residuals $\rb = \m - \fb(\hat\xib)$. Let $\delta\xib$ be the small perturbation that moves the true pose $\xib$ to the estimate $\hat\xib$. Linearizing $\fb$ around the true pose, $\fb(\hat\xib) \approx \fb(\xib) + \J \, \delta\xib$, gives the residuals
\begin{equation}
  \rb \approx \epsb - \J \, \delta\xib . \label{eq:residuals_linear}
\end{equation}
At the minimum, the gradient of the cost, $\partial E / \partial \delta\xib = -\J^\tran \rb$, vanishes. $\delta\xib$ therefore solves the normal equations $\J^\tran \J \, \delta\xib = \J^\tran \epsb$, the linear least-squares fit of a camera motion to the noise,
\begin{equation}
  \delta\xib = (\J^\tran \J)^{-1} \J^\tran \epsb . \label{eq:pose_error}
\end{equation}
Its covariance is
\[
  \inCs \Sb_\C = \mathbb{E}[\delta\xib \, \delta\xib^\tran]
  = (\J^\tran \J)^{-1} \J^\tran \, \mathbb{E}[\epsb \epsb^\tran] \, \J (\J^\tran \J)^{-1}
  = \sigma^2 (\J^\tran \J)^{-1},
\]
OpenCV orders the Jacobian columns as (rotation, translation). We reorder them to (translation, rotation), so that $\inCs \Sb_\C$ has the layout of a ROS pose covariance.

\paragraph{Estimating the pixel noise $\sigma$.}
$\sigma^2$ is the variance of a single pixel coordinate, so it is estimated from the $2N$ coordinate residuals $\rb$. Substituting \eqref{eq:pose_error} into \eqref{eq:residuals_linear} gives
\[
  \rb = (\Ib_{2N} - \Pb)\, \epsb, \qquad \Pb = \J (\J^\tran \J)^{-1} \J^\tran .
\]
The column space of $\J$ holds the corner displacements $\J \, \delta\xib$ that a small camera motion produces. $\Pb$ is the orthogonal projection onto it: $\Pb \epsb$ is the part of the noise that a camera motion explains, and the residuals are the rest, orthogonal to every camera motion.

$\Pb$ is symmetric and idempotent, $\Pb^2 = \J (\J^\tran \J)^{-1} \J^\tran \J (\J^\tran \J)^{-1} \J^\tran = \Pb$, so $\Ib_{2N} - \Pb$ is too,
\begin{equation}
  (\Ib_{2N} - \Pb)^\tran (\Ib_{2N} - \Pb) = (\Ib_{2N} - \Pb)^2 = \Ib_{2N} - 2\Pb + \Pb^2 = \Ib_{2N} - \Pb . \label{eq:idempotent}
\end{equation}
For any matrix $\bm{M}$, the noise coordinates are uncorrelated, $\mathbb{E}[\varepsilon_i \varepsilon_j] = 0$ for $i \neq j$, and $\mathbb{E}[\varepsilon_i^2] = \sigma^2$, so only the diagonal of $\bm{M}$ remains,
\begin{equation}
  \mathbb{E}[\epsb^\tran \bm{M} \epsb]
  = \sum_{i,j} M_{ij} \, \mathbb{E}[\varepsilon_i \varepsilon_j]
  = \sigma^2 \sum_i M_{ii}
  = \sigma^2 \operatorname{tr}(\bm{M}) . \label{eq:trace}
\end{equation}
The expected residual sum of squares follows from \eqref{eq:idempotent}, \eqref{eq:trace} and the cyclic property of the trace, $\operatorname{tr}(\bm{A}\bm{B}) = \operatorname{tr}(\bm{B}\bm{A})$,
\begin{equation}
  \begin{aligned}
    \mathbb{E}[\rb^\tran \rb]
     & = \mathbb{E}[\epsb^\tran (\Ib_{2N} - \Pb)^\tran (\Ib_{2N} - \Pb) \epsb]
    \stackrel{\eqref{eq:idempotent}}{=} \mathbb{E}[\epsb^\tran (\Ib_{2N} - \Pb) \epsb]
    \stackrel{\eqref{eq:trace}}{=} \sigma^2 \operatorname{tr}(\Ib_{2N} - \Pb)                        \\
     & = \sigma^2 \left( 2N - \operatorname{tr}\left( (\J^\tran \J)^{-1} \J^\tran \J \right) \right)
    = (2N - 6)\, \sigma^2 .
  \end{aligned}
  \label{eq:rss}
\end{equation}
Dividing the residual sum of squares by $2N - 6$ therefore gives an unbiased estimate of $\sigma^2$, $\mathbb{E}[\hat\sigma^2] = \sigma^2$~\cite{statproofbook_resvar_biasp}, with
\begin{equation}
  \hat\sigma^2 = \frac{\rb^\tran \rb}{2N - 6} = \frac{1}{2N - 6} \sum_{i=1}^{N} \left( \Delta u_i^2 + \Delta v_i^2 \right), \label{eq:sigma_hat}
\end{equation}
which needs $N \geq 4$ corners. A single marker always provides 4.

The $6$ in \eqref{eq:rss} is the trace of $\Pb$, the number of pose parameters. Fitting them makes the residuals satisfy the 6 normal equations $\J^\tran \rb = \Zb$, so only $2N - 6$ of them vary freely. Fitting a mean to $n$ samples is the case $\J = [1, \dots, 1]^\tran$: $\Pb$ replaces every sample with the mean, $\operatorname{tr}(\Pb) = 1$, the normal equation states that the residuals sum to zero, and \eqref{eq:sigma_hat} becomes the sample variance with its $n - 1$.

\paragraph{Lower bound.}
With one marker, $2N - 6 = 2$: the normal equations leave the 8 residuals only 2 independent components, so $\hat\sigma^2$ is (in effect) an average of just two squared noise samples. Two samples are too few to average out (in about 10\,\% of the frames, $\hat\sigma^2$ is below a tenth of the true variance\footnote{To first order, \eqref{eq:residuals_linear} is a linear regression model with design matrix $\J$, so $\rb^\tran \rb / \sigma^2 = (2N - 6)\,\hat\sigma^2 / \sigma^2$ follows a $\chi^2$ distribution with $2N - 6$ degrees of freedom~\cite{statproofbook_rssdist}. For one marker, that is $\chi^2_2$, the exponential distribution with mean 2~\cite{wikipedia_chi2}, so $P(\hat\sigma^2 < k \sigma^2) = P(\chi^2_2 < 2k) = 1 - e^{-k}$.}). A localization filter would trust such a near-zero covariance blindly. The pixel noise should therefore be floored at the parameter $\sigma_\mathrm{min}$,
\[
  \sigma^2 = \max\left(\hat\sigma^2, \sigma_\mathrm{min}^2\right),
  \qquad
  \inCs \Sb_\C \approx \sigma^2 (\J^\tran \J)^{-1}.
\]
The noise of a static camera's corner detections, or the RMS error of the camera calibration, gives a value for $\sigma_\mathrm{min}$.
The residuals capture the errors that make the markers inconsistent with each other, such as corner bias, a mis-surveyed marker or calibration errors, and they change with the blur, lighting and distance in each frame. Errors that shift all markers consistently, such as an error in $\inB \Hb_\C$, leave no residual and are not modeled.

\begin{thebibliography}{9}
  \bibitem{statproofbook_rssdist}
  The Book of Statistical Proofs, ``Distribution of residual sum of squares in multiple linear regression'', \url{https://statproofbook.github.io/P/mlr-rssdist}.
  \bibitem{statproofbook_resvar_biasp}
  The Book of Statistical Proofs, ``Maximum likelihood estimator of variance in multiple linear regression is biased'', \url{https://statproofbook.github.io/P/resvar-biasp}.
  \bibitem{wikipedia_chi2}
  Wikipedia, ``Chi-squared distribution'', section ``Related distributions'', \url{https://en.wikipedia.org/wiki/Chi-squared_distribution}.
\end{thebibliography}
\end{document}
